\subsection{Dimension-block obstructions and their scope}
\label{subsec:dimension-block}

\begin{proposition}\label{prop:pivotal-dimension-block}
The rings $R7\_3\_695$, $R7\_3\_783$, and $R7\_3\_1015$ have no
pivotal categorification over $\mathbb C$.
\end{proposition}
\begin{proof}
For a spherical fusion category whose global dimension $D$ generates
a real quadratic field of discriminant $f$, the nontrivial-character
case of Schopieray's dimension-block theorem
\cite[Theorem~5.1]{SchopierayDimensionBlock} gives
\[
 \left(\frac1D-\frac1{D'}\right)^2\leq\frac1{e^2 f}
 \leq\frac1f,
\]
where $e$ is a positive integer. The relevant source dimensions and
exact values are
\[
\begin{array}{c|c|r|c}
 \text{ring}&D&f&(1/D-1/D')^2\\ \hline
 R7\_3\_695&420+70\sqrt{35}&140&1/35\\
 R7\_3\_783, R7\_3\_1015&120+30\sqrt{15}&60&1/15.
\end{array}
\]
Both conjugates violate the necessary inequality in each case.

It remains to exclude other dimension characters. Their global
dimensions must be eigenvalues of the regular Casimir. For
$R7\_3\_695$, the other roots are $5,7$, and
$\operatorname{Norm}(D)=4900$; the norms of $5/D$ and $7/D$ are
$1/196,1/100$. For the other two rings the rational roots are $6,15$,
and $\operatorname{Norm}(D)=900$; the corresponding norms are
$1/25,1/4$. Thus none of the rational alternatives satisfies
formal-codegree divisibility
\cite[Corollary~2.14]{OstrikRank3}. In the noncommutative case it is
harmless to include the Casimir root from the matrix block in this
necessary superset.

Finally every pivotal dimension character for these rings is
spherical. All objects of $R7\_3\_695$ are self-dual. In the other
two rings, the three invertibles form $C_3$, a self-dual singleton
$E$ satisfies $gE=E$, and the other three simples form one $C_3$-orbit.
Nonvanishing of $\dim(E)$ forces $\dim(g)=1$; the three orbit dimensions
are equal, so every simple and its dual have the same dimension.
\end{proof}

This proposition alone is a pivotal obstruction, not a general one.
The full sphericalization doubles the global dimension, dividing the
left-hand side of the displayed inequality by four; the strict
violations then become equalities. No general exclusion is inferred
solely from this dimension-block argument. The companion verifier
keeps this scope explicit and checks all numerical and tensor
hypotheses exactly.
